\documentclass[12pt,a4paper]{article}
\usepackage{bendingforks}

% =====================================================================
%  PREVENTIVO DEI COSTI E ASSUNZIONE DEGLI IMPEGNI - documentazione esterna
%  Vincoli: costo totale >= 12.000 €, 80-95 ore produttive per persona
%  (7 persone: da 560 a 665 ore totali). Tutti devono ricoprire tutti i ruoli.
%  I numeri in rosso sono la distribuzione indicativa delle slide:
%  vanno sostituiti con quella decisa dal gruppo. Date AAAA-MM-GG.
% =====================================================================
\bfsetup{Preventivo dei Costi e Assunzione degli Impegni}{Capitolato C3 -- Fossdot}{0.1.0}{Documento}{Esterno}
\bfredattori{\ph{Nome Cognome}}
\bfverificatori{\ph{Nome Cognome}}
\bfdestinatari{\emph{Prof. Tullio Vardanega} \\ & \emph{Prof. Riccardo Cardin}}
\bfabstract{Preventivo economico e orario del gruppo \emph{Bending Forks} per la realizzazione del capitolato C3 \emph{Fossdot} proposto da \emph{Uranion}: ripartizione delle ore per ruolo e per membro, costo totale, analisi delle disponibilità e assunzione formale degli impegni.}

\begin{document}
\bfcopertina
\bfindice

\begin{bfregistro}
0.1.0 & \ph{2026-10-27} & \ph{Nome Cognome} & Creazione della struttura del documento \bfsep
\end{bfregistro}

% =====================================================================
\section{Introduzione}

\subsection{Capitolato selezionato}
Il presente documento riporta il preventivo dei costi e gli impegni assunti dal gruppo \emph{Bending Forks} per il capitolato C3 \emph{Fossdot}, proposto dall'azienda \emph{Uranion}.

\subsection{Periodo di progetto stimato}
\begin{itemize}
\tightlist
\item \textbf{Inizio:} \ph{AAAA-MM-GG} (aggiudicazione prevista il 2026-11-03)
\item \textbf{Revisione RTB:} \ph{AAAA-MM-GG}
\item \textbf{Revisione PB:} \ph{AAAA-MM-GG}
\item \textbf{Consegna finale prevista:} \ph{AAAA-MM-GG}
\end{itemize}

% =====================================================================
\section{Ripartizione dei ruoli preventivata}

\subsection{Distribuzione delle ore per ruolo}

\begin{center}
\begin{tabularx}{\linewidth}{>{\hsize=1.3\hsize\raggedright\arraybackslash}X>{\hsize=0.9\hsize\centering\arraybackslash}X>{\hsize=0.9\hsize\centering\arraybackslash}X>{\hsize=0.9\hsize\centering\arraybackslash}X}
\bfheadrow\bfh{Ruolo} & \bfh{Ore} & \bfh{Costo orario} & \bfh{Costo} \bfsep
Responsabile & \ph{60} & 30\,€ & \ph{1.800\,€} \bfsep
Amministratore & \ph{70} & 20\,€ & \ph{1.400\,€} \bfsep
Analista & \ph{100} & 25\,€ & \ph{2.500\,€} \bfsep
Progettista & \ph{120} & 25\,€ & \ph{3.000\,€} \bfsep
Programmatore & \ph{180} & 15\,€ & \ph{2.700\,€} \bfsep
Verificatore & \ph{65} & 15\,€ & \ph{975\,€} \bfsep
\textbf{Totale} & \textbf{\ph{595}} & -- & \textbf{\ph{12.375\,€}} \bfsep
\end{tabularx}
\end{center}

\ph{Commento: motivare il peso dei ruoli (ad es. più ore di Programmatore e Verificatore per la parte hardware e firmware).}

\subsection{Ore per membro e per ruolo}
Ogni membro ricopre tutti i ruoli nel corso del progetto. Sigle: Re = Responsabile, Am = Amministratore, An = Analista, Pt = Progettista, Pr = Programmatore, Ve = Verificatore.

\begin{center}
\begin{tabularx}{\linewidth}{>{\hsize=1.9\hsize\raggedright\arraybackslash}X*{6}{>{\hsize=0.65\hsize\centering\arraybackslash}X}>{\hsize=0.75\hsize\centering\arraybackslash}X}
\bfheadrow\bfh{Membro} & \bfh{Re} & \bfh{Am} & \bfh{An} & \bfh{Pt} & \bfh{Pr} & \bfh{Ve} & \bfh{Totale} \bfsep
Marco Beltrame & \ph{} & \ph{} & \ph{} & \ph{} & \ph{} & \ph{} & \ph{85} \bfsep
Lorenzo Milanese & \ph{} & \ph{} & \ph{} & \ph{} & \ph{} & \ph{} & \ph{85} \bfsep
Filippo Rampazzo & \ph{} & \ph{} & \ph{} & \ph{} & \ph{} & \ph{} & \ph{85} \bfsep
Nicholas Salmaso & \ph{} & \ph{} & \ph{} & \ph{} & \ph{} & \ph{} & \ph{85} \bfsep
Luca Stangherlin & \ph{} & \ph{} & \ph{} & \ph{} & \ph{} & \ph{} & \ph{85} \bfsep
Aldo Toffoli & \ph{} & \ph{} & \ph{} & \ph{} & \ph{} & \ph{} & \ph{85} \bfsep
Matteo Wu & \ph{} & \ph{} & \ph{} & \ph{} & \ph{} & \ph{} & \ph{85} \bfsep
\textbf{Totale} & \ph{60} & \ph{70} & \ph{100} & \ph{120} & \ph{180} & \ph{65} & \textbf{\ph{595}} \bfsep
\end{tabularx}
\end{center}

\subsection{Rotazione dei ruoli pianificata}
\ph{Ogni quanto ruotano i ruoli (ad es. a ogni sprint di due settimane), con quale criterio e come si garantisce che redattore e verificatore di un documento siano persone diverse.}

\subsection{Copertura dei ruoli}
\ph{Conferma che ogni membro ricopre ogni ruolo almeno una volta, con rimando alla tabella precedente.}

% =====================================================================
\section{Preventivo economico}

\subsection{Costo per ruolo}
Vedi la tabella della sezione 2.1.

\subsection{Costo per periodo}

\begin{center}
\begin{tabularx}{\linewidth}{>{\hsize=1.4\hsize\raggedright\arraybackslash}X>{\hsize=0.8\hsize\centering\arraybackslash}X>{\hsize=0.8\hsize\centering\arraybackslash}X}
\bfheadrow\bfh{Periodo} & \bfh{Ore} & \bfh{Costo} \bfsep
Fino alla RTB & \ph{ore} & \ph{€} \bfsep
Dalla RTB alla PB & \ph{ore} & \ph{€} \bfsep
\textbf{Totale} & \textbf{\ph{595}} & \textbf{\ph{12.375\,€}} \bfsep
\end{tabularx}
\end{center}

\subsection{Costo totale preventivato}
Il costo totale preventivato è di \textbf{\ph{12.375\,€}} per \textbf{\ph{595}} ore produttive. Questo importo costituisce il limite superiore invalicabile del costo del progetto.

% =====================================================================
\section{Analisi delle disponibilità}

\subsection{Impegni accademici dei membri}
\ph{Esami e sessioni, lezioni con obbligo di frequenza, ecc.}

\subsection{Altri impegni rilevanti}
\ph{Lavoro, stage, altri impegni personali che riducono la disponibilità.}

\subsection{Calendario condiviso delle disponibilità}
\ph{Dove si trova il calendario (ad es. Google Calendar del gruppo) e come viene aggiornato.}

\subsection{Periodi critici}

\begin{center}
\begin{tabularx}{\linewidth}{>{\hsize=0.9\hsize\raggedright\arraybackslash}X>{\hsize=1.4\hsize\raggedright\arraybackslash}X>{\hsize=0.7\hsize\raggedright\arraybackslash}X}
\bfheadrow\bfh{Periodo} & \bfh{Motivo} & \bfh{Membri coinvolti} \bfsep
\ph{AAAA-MM-GG / AAAA-MM-GG} & \ph{Sessione invernale d'esami} & \ph{Tutti} \bfsep
\ph{AAAA-MM-GG / AAAA-MM-GG} & \ph{Vacanze di Natale} & \ph{Tutti} \bfsep
\end{tabularx}
\end{center}

% =====================================================================
\section{Assunzione formale degli impegni}

\subsection{Ore garantite per membro}
Ciascun membro del gruppo \emph{Bending Forks} si impegna a garantire il seguente numero di ore produttive.

\begin{center}
\begin{tabularx}{\linewidth}{>{\hsize=1.1\hsize\raggedright\arraybackslash}X>{\hsize=0.6\hsize\centering\arraybackslash}X>{\hsize=0.6\hsize\centering\arraybackslash}X>{\hsize=0.7\hsize\centering\arraybackslash}X}
\bfheadrow\bfh{Membro} & \bfh{Matricola} & \bfh{Ore garantite} & \bfh{Firma} \bfsep
Marco Beltrame & 2111021 & \ph{85} & \bfsep
Lorenzo Milanese & 2117110 & \ph{85} & \bfsep
Filippo Rampazzo & 2101097 & \ph{85} & \bfsep
Nicholas Salmaso & 2080038 & \ph{85} & \bfsep
Luca Stangherlin & 2137996 & \ph{85} & \bfsep
Aldo Toffoli & 2113202 & \ph{85} & \bfsep
Matteo Wu & 2075514 & \ph{85} & \bfsep
\end{tabularx}
\end{center}

\subsection{Scadenza ultima di consegna prevista}
Il gruppo si impegna a completare il progetto entro il \textbf{\ph{AAAA-MM-GG}}.

\subsection{Firma del responsabile}
\vspace{20pt}
\noindent\bffirma{Il Responsabile del gruppo:}{\ph{Nome Cognome}}

\end{document}
